\subsection{prime-log 临界势的 coboundary 化与 Abel--Mellin 边界留数刚性}\label{subsec:conclusion-prime-log-coboundary-abel-mellin-rigidity}
沿用定理 \ref{thm:conclusion-prime-log-calibration-zero-lyapunov-product}、定理 \ref{thm:xi-projective-pressure-path-holder-convexity}、定理 \ref{thm:xi-projective-pressure-integer-spectral-certificate}、定理 \ref{thm:xi-projective-moment-tower-anchors-normalized-pressure}、定理 \ref{thm:xi-debranges-canonical-extreme-zk} 与猜想 \ref{conj:conclusion-prime-log-canonical-system} 的记号。上一小节已经说明：prime-log 校准给出零 Lyapunov 的临界归一化、加权 prime-end 的对数级边界流，以及 RH 候选算子所不可回避的可数无限秩 prime-register 本体。若再把这条校准链与实参数射影压力、Abel--Mellin 提升及 de Branges 的边界参数化并读，则 centered prime-log 势还可进一步压缩为如下五条后果：差额序列把临界势精确写成一维 coboundary；该 coboundary 对 \(\mathcal L_q\) 只能造成次指数级的 boundary twist，因而不改变 bulk 谱压；时间指标上的 Abel--Mellin 边界流在 \(\Re s>0\) 中只保留一条留数为 \(2\alpha\log B\) 的单极点；维数参数 \(\alpha\) 的全部非全纯依赖都压缩到 \(\mathbb C\cdot \zeta(s)\) 这一维方向；因此若 RH 的 canonical-system 路线能够闭合，则 \(\alpha\) 只能作为边界质量而非 bulk Hamiltonian 形变出现。

\begin{theorem}[prime-log 临界势的精确 coboundary 化]\label{thm:conclusion-prime-log-exact-coboundary}
沿用定理 \ref{thm:conclusion-prime-log-calibration-zero-lyapunov-product} 的记号，并定义 centered prime-log 势
\[
\psi_\alpha(j):=y_jw_j-c
\qquad (j\ge 1).
\]
则对每个 \(j\ge 1\) 都有精确恒等式
\[
\psi_\alpha(j)=\delta_{j-1}-\delta_j.
\]
因而
\[
e^{\psi_\alpha(j)}=\frac{e^{\delta_{j-1}}}{e^{\delta_j}},
\qquad
\exp\!\Bigl(\sum_{j=1}^{n}\psi_\alpha(j)\Bigr)=e^{-\delta_n}
\qquad (n\ge 1).
\]
特别地，
\[
0\le \delta_n<w_{n+1}=\log p_{m_0+2n+2},
\qquad
p_{m_0+2n+2}^{-1}<e^{-\delta_n}\le 1.
\]
\end{theorem}
\begin{proof}
由定理 \ref{thm:conclusion-prime-log-calibration-zero-lyapunov-product} 的递归关系，
\[
\delta_j-\delta_{j-1}
=
\bigl(\delta_{j-1}+c-y_jw_j\bigr)-\delta_{j-1}
=
c-y_jw_j.
\]
移项即得
\[
\psi_\alpha(j)=y_jw_j-c=\delta_{j-1}-\delta_j.
\]
逐项指数化便有
\[
e^{\psi_\alpha(j)}=\frac{e^{\delta_{j-1}}}{e^{\delta_j}}.
\]
再对 \(j=1,\dots,n\) 连乘并用 \(\delta_0=0\)，得到
\[
\exp\!\Bigl(\sum_{j=1}^{n}\psi_\alpha(j)\Bigr)
=
\prod_{j=1}^{n}\frac{e^{\delta_{j-1}}}{e^{\delta_j}}
=
e^{-\delta_n}.
\]
最后，由定理 \ref{thm:conclusion-prime-log-calibration-zero-lyapunov-product} 已知
\(
0\le \delta_n<w_{n+1}
\)，指数化即得所示界。证毕。
\end{proof}

\begin{theorem}[prime-log 边界扭曲不改变射影算子的 bulk 谱压]\label{thm:conclusion-prime-log-boundary-twist-pressure-invariance}
沿用定义 \ref{def:pom-projective-transfer-operator} 与定理 \ref{thm:conclusion-prime-log-exact-coboundary} 的记号。对 \(q\ge 1\) 与 \(b\in A_{\mathrm{out}}\)，定义分支算子
\[
(\mathcal L_{q,b}f)(w):=g_b(w)^q f(\Phi_b(w)),
\qquad
\mathcal L_q=\sum_{b\in A_{\mathrm{out}}}\mathcal L_{q,b}.
\]
对任意输出词 \(\mathbf b=b_1\cdots b_n\in A_{\mathrm{out}}^n\)，记
\[
\mathcal L_{q,\mathbf b}:=\mathcal L_{q,b_n}\cdots \mathcal L_{q,b_1},
\qquad
\mathcal L_{q,\mathbf b}^{(\alpha)}
:=
\exp\!\Bigl(\sum_{j=1}^{n}\psi_\alpha(j)\Bigr)\mathcal L_{q,\mathbf b}.
\]
再定义对应的 \(n\)-步校准扭曲算子
\[
\mathcal L_q^{(\alpha,n)}
:=
\sum_{|\mathbf b|=n}\mathcal L_{q,\mathbf b}^{(\alpha)}.
\]
则对每个 \(n\ge 1\) 与每个 \(|\mathbf b|=n\)，都有
\[
\mathcal L_{q,\mathbf b}^{(\alpha)}=e^{-\delta_n}\mathcal L_{q,\mathbf b},
\qquad
\mathcal L_q^{(\alpha,n)}=e^{-\delta_n}\mathcal L_q^n.
\]
因而对任意一致算子范数，
\[
e^{-w_{n+1}}\|\mathcal L_{q,\mathbf b}\|
\le
\|\mathcal L_{q,\mathbf b}^{(\alpha)}\|
\le
\|\mathcal L_{q,\mathbf b}\|,
\]
\[
e^{-w_{n+1}}\|\mathcal L_q^n\|
\le
\|\mathcal L_q^{(\alpha,n)}\|
\le
\|\mathcal L_q^n\|.
\]
特别地，
\[
\lim_{n\to\infty}\frac1n\log \|\mathcal L_q^{(\alpha,n)}\|
=
\log \rho(\mathcal L_q).
\]
并且对任意固定路径族 \(\mathbf b^{(n)}\in A_{\mathrm{out}}^n\)，若
\[
\lim_{n\to\infty}\frac1n\log \|\mathcal L_{q,\mathbf b^{(n)}}\|
\]
存在，则
\[
\lim_{n\to\infty}\frac1n\log \|\mathcal L_{q,\mathbf b^{(n)}}^{(\alpha)}\|
=
\lim_{n\to\infty}\frac1n\log \|\mathcal L_{q,\mathbf b^{(n)}}\|.
\]
换言之，prime-log 校准只能以 boundary twist 的方式进入 \(\mathcal L_q\)，而不能改变其 bulk 谱压；正文中的连续归一化压力 \(\Lambda(t)\) 因而也不产生 bulk 位移。
\end{theorem}
\begin{proof}
由定理 \ref{thm:conclusion-prime-log-exact-coboundary}，
\[
\sum_{j=1}^{n}\psi_\alpha(j)=-\delta_n.
\]
故对每个 \(|\mathbf b|=n\)，都有
\[
\mathcal L_{q,\mathbf b}^{(\alpha)}=e^{-\delta_n}\mathcal L_{q,\mathbf b}.
\]
对所有长度为 \(n\) 的词求和，便得
\[
\mathcal L_q^{(\alpha,n)}
=
\sum_{|\mathbf b|=n}e^{-\delta_n}\mathcal L_{q,\mathbf b}
=
e^{-\delta_n}\sum_{|\mathbf b|=n}\mathcal L_{q,\mathbf b}
=
e^{-\delta_n}\mathcal L_q^n.
\]
由定理 \ref{thm:conclusion-prime-log-calibration-zero-lyapunov-product} 的界
\(
0\le \delta_n<w_{n+1}
\)，立得
\[
e^{-w_{n+1}}
<
e^{-\delta_n}
\le 1,
\]
从而得到两组范数夹逼。

对全局 \(n\)-步扭曲算子，
\[
\frac1n\log \|\mathcal L_q^{(\alpha,n)}\|
=
\frac1n\log \|\mathcal L_q^n\|
-\frac{\delta_n}{n}.
\]
由于 \(w_{n+1}=\log p_{m_0+2n+2}=o(n)\) 且 \(0\le \delta_n<w_{n+1}\)，故
\[
\frac{\delta_n}{n}\longrightarrow 0.
\]
再由 Gelfand 谱半径公式，
\[
\lim_{n\to\infty}\frac1n\log \|\mathcal L_q^n\|
=
\log\rho(\mathcal L_q),
\]
遂得
\[
\lim_{n\to\infty}\frac1n\log \|\mathcal L_q^{(\alpha,n)}\|
=
\log\rho(\mathcal L_q).
\]

同理，对任意路径族 \(\mathbf b^{(n)}\in A_{\mathrm{out}}^n\)，都有
\[
\frac1n\log \|\mathcal L_{q,\mathbf b^{(n)}}^{(\alpha)}\|
=
\frac1n\log \|\mathcal L_{q,\mathbf b^{(n)}}\|
-\frac{\delta_n}{n},
\]
故一旦右端极限存在，左端极限便存在且相等。最后，正文中的连续归一化压力 \(\Lambda(t)\) 只是在 Perron 谱压上再减去常数 \(\log|A_{\mathrm{out}}|\)（见定理 \ref{thm:xi-projective-moment-tower-anchors-normalized-pressure}），因而同样不受该 boundary twist 的 bulk 修正。证毕。
\end{proof}

\begin{theorem}[Abel--Mellin 边界流的单极点刚性与全纯余项]\label{thm:conclusion-prime-log-abel-mellin-pole-rigidity}
沿用前述记号，定义时间指标上的 Abel--Mellin 边界流
\[
\mathcal D_\alpha(s):=
\sum_{j=1}^{\infty}\frac{y_jw_j}{j^s}
\qquad (\Re s>1),
\]
并在同一半平面上定义 centered 部分
\[
\mathcal J_\alpha(s):=
\sum_{j=1}^{\infty}\frac{\psi_\alpha(j)}{j^s}.
\]
则 \(\mathcal J_\alpha\) 在 \(\Re s>0\) 上存在唯一的全纯延拓，并且该延拓满足绝对收敛表示
\[
\mathcal J_\alpha(s)
=
\sum_{j=1}^{\infty}
\delta_j\bigl((j+1)^{-s}-j^{-s}\bigr)
\qquad (\Re s>0).
\]
在原定义半平面 \(\Re s>1\) 上有严格分解
\[
\mathcal D_\alpha(s)=c\,\zeta(s)+\mathcal J_\alpha(s).
\]
因此 \(\mathcal D_\alpha\) 在 \(\Re s>0\) 中唯一可能的奇性是 \(s=1\) 处的简单极点，并且
\[
\operatorname*{Res}_{s=1}\mathcal D_\alpha(s)=c=2\alpha\log B.
\]
\end{theorem}
\begin{proof}
在 \(\Re s>1\) 上，由定理 \ref{thm:conclusion-prime-log-exact-coboundary}，
\[
\mathcal J_\alpha(s)
=
\sum_{j=1}^{\infty}\frac{\delta_{j-1}-\delta_j}{j^s}.
\]
由于 \(\delta_0=0\)，把第一项移位可得
\[
\sum_{j=1}^{\infty}\frac{\delta_{j-1}}{j^s}
=
\sum_{j=1}^{\infty}\frac{\delta_j}{(j+1)^s},
\]
从而
\[
\mathcal J_\alpha(s)
=
\sum_{j=1}^{\infty}
\delta_j\bigl((j+1)^{-s}-j^{-s}\bigr).
\]
再由定理 \ref{thm:conclusion-prime-log-calibration-zero-lyapunov-product}，
\[
0\le \delta_j<w_{j+1}=\log p_{m_0+2j+2}\ll \log(j+2).
\]
对任意紧集 \(K\subset\{\,\Re s>0\,\}\)，存在 \(\sigma_0>0\) 使得 \(\Re s\ge \sigma_0\) 对所有 \(s\in K\) 成立；而
\[
(j+1)^{-s}-j^{-s}=O_K(j^{-\sigma_0-1}).
\]
故
\[
\delta_j\bigl((j+1)^{-s}-j^{-s}\bigr)
=
O_K\!\bigl(\log(j+2)\,j^{-\sigma_0-1}\bigr),
\]
其右端可求和。由 Weierstrass 判别法，上述级数在每个这样的紧集上一致绝对收敛，从而给出 \(\mathcal J_\alpha\) 在 \(\Re s>0\) 的唯一全纯延拓。

另一方面，在 \(\Re s>1\) 上有
\[
y_jw_j=c+\psi_\alpha(j),
\]
逐项求和即得
\[
\mathcal D_\alpha(s)
=
\sum_{j=1}^{\infty}\frac{c}{j^s}
+\sum_{j=1}^{\infty}\frac{\psi_\alpha(j)}{j^s}
=
c\,\zeta(s)+\mathcal J_\alpha(s).
\]
由于 \(\zeta(s)\) 在 \(\Re s>0\) 中仅于 \(s=1\) 具有留数为 \(1\) 的简单极点，而 \(\mathcal J_\alpha\) 在该半平面全纯，故 \(\mathcal D_\alpha\) 的唯一可能奇性正是 \(s=1\) 处的简单极点，其留数等于 \(c=2\alpha\log B\)。证毕。
\end{proof}

\begin{corollary}[维数参数的非全纯依赖只有一维边界方向]\label{cor:conclusion-prime-log-alpha-singular-rank-one}
\leanverified{paper\_conclusion\_prime\_log\_alpha\_singular\_rank\_one}
设 \(0<\alpha_1\neq \alpha_2<\infty\)。对 \(\nu=1,2\)，记
\[
c_\nu:=2\alpha_\nu\log B,
\]
并令 \(y_j^{(\nu)}\)、\(\delta_j^{(\nu)}\)、\(\psi_{\alpha_\nu}(j)\)、\(\mathcal D_{\alpha_\nu}(s)\) 为相应的 prime-log 校准量与 Abel--Mellin 边界流。则在 \(\Re s>1\) 上有
\[
\mathcal D_{\alpha_1}(s)-\mathcal D_{\alpha_2}(s)
=
2(\alpha_1-\alpha_2)\log B\,\zeta(s)+H_{12}(s),
\]
其中 \(H_{12}\) 在 \(\Re s>0\) 上全纯。特别地，\(\alpha\)-依赖的全部非全纯部分都落在
\[
\mathbb C\cdot \zeta(s)
\]
这一维方向上。
\end{corollary}
\begin{proof}
对 \(\nu=1,2\)，由定理 \ref{thm:conclusion-prime-log-abel-mellin-pole-rigidity}，
\[
\mathcal D_{\alpha_\nu}(s)
=
c_\nu\zeta(s)+\mathcal J_{\alpha_\nu}(s),
\]
其中 \(\mathcal J_{\alpha_\nu}\) 在 \(\Re s>0\) 上全纯。两式相减便得
\[
\mathcal D_{\alpha_1}(s)-\mathcal D_{\alpha_2}(s)
=
(c_1-c_2)\zeta(s)
+\bigl(\mathcal J_{\alpha_1}(s)-\mathcal J_{\alpha_2}(s)\bigr).
\]
记
\[
H_{12}(s):=\mathcal J_{\alpha_1}(s)-\mathcal J_{\alpha_2}(s),
\]
则 \(H_{12}\) 在 \(\Re s>0\) 上全纯；再把 \(c_1-c_2\) 写成 \(2(\alpha_1-\alpha_2)\log B\) 即得所示公式。证毕。
\end{proof}

结合定理 \ref{thm:xi-debranges-canonical-extreme-zk} 的 canonical-system 边界参数化可知，若要把 \(\alpha\) 接入 RH 的谱实现，则它只能表现为边界质量、Weyl \(m\)-函数的 Herglotz 规一化常数或等价的谱移参数；若把 \(\alpha\) 作为 bulk Hamiltonian 的独立形变引入，便会制造超出 \(\mathbb C\cdot\zeta(s)\) 的高维奇性家族，从而与上式的一维留数刚性矛盾。

\begin{conjecture}[prime-log boundary-mass canonical system]\label{conj:conclusion-prime-log-boundary-mass-canonical-system}
存在一族定义在可数 prime-register 轴上的 canonical system
\[
\mathcal H_\alpha,
\]
以及与之对应的 Abel--Mellin 边界散射核 \(\mathcal K_\alpha(s)\)，使其正则化 Fredholm 行列式满足
\[
-\frac{\mathrm d}{\mathrm ds}\log \det_{\mathrm{ren}}(I-\mathcal K_\alpha(s))
=
\mathcal D_\alpha(s)
=
2\alpha\log B\,\zeta(s)+\mathcal J_\alpha(s),
\]
其中 \(\mathcal J_\alpha\) 为定理 \ref{thm:conclusion-prime-log-abel-mellin-pole-rigidity} 所给出的 \(\Re s>0\) 上全纯余项。并且该族满足：
\begin{enumerate}
  \item bulk Hamiltonian 与 \(\alpha\) 无关；
  \item \(\alpha\) 仅通过一个边界质量参数进入 Weyl \(m\)-函数；
  \item RH 若成立，则等价于该边界质量族在临界线上的自伴谱移保持纯虚性。
\end{enumerate}
\end{conjecture}

该猜想可视为猜想 \ref{conj:conclusion-prime-log-canonical-system} 的边界质量强化版。定理 \ref{thm:conclusion-prime-log-boundary-twist-pressure-invariance} 已排除 \(\alpha\) 通过 bulk 谱压进入；定理 \ref{thm:conclusion-prime-log-abel-mellin-pole-rigidity} 与推论 \ref{cor:conclusion-prime-log-alpha-singular-rank-one} 又把全部非全纯信息压缩到 \(\zeta\)-型单极点的一维方向；再结合定理 \ref{thm:conclusion-prime-log-rh-operator-countable-prime-rank} 的可数无限秩障碍与定理 \ref{thm:xi-debranges-canonical-extreme-zk} 的边界参数化，\(\alpha\) 的可实现位置便只剩边界质量/谱移这一类一维缺陷通道。

\noindent
由此，prime-log calibration 在结论层形成如下刚性链：
\[
\text{centered prime-log potential}
\Longrightarrow
\text{exact coboundary}
\Longrightarrow
\text{subexponential boundary twist of }\mathcal L_q
\Longrightarrow
\text{one-dimensional Abel--Mellin singular current}.
\]
因此，与 de Branges/canonical system 接口兼容的 RH 候选算子不再允许把 \(\alpha\) 作为 bulk 形变参数引入；它若存在，只能是一个 bulk 谱压保持不变、边界留数一维可调、并在 Abel--Mellin 提升后恢复 \(\zeta\)-型单极点的可数无限秩边界质量模型。
\endinput
